\documentclass[11pt]{article}
\usepackage{mathblatt}
% Präambel des Setzers (werkzeuge/setzer-praeambel.tex, wortgleich); dazu \gteil (Muster T2U)
\makeatletter
\setlength{\mbpfbe}{0mm}% keine Punkte (pruefung.md § 5: nur beim Prüfstein)
\newlength{\gEin}\setlength{\gEin}{0pt}
\let\gSchrift\relax
\renewcommand{\pfkreuzzeile}[1]{\par\noindent\hangindent3.2em\hangafter1\hspace*{1.6em}$\square$\ #1\par}
\newcommand{\gkreuz}[1]{$\square$\ #1\hspace{2em}}
\newcommand{\gkreuzzeile}[2]{\par\noindent\hangindent2.2em\hangafter1\mbox{#1\,$\square$}\ #2\par}% Buchstabe, Option: eine je Zeile (neben Skizze)
% links, was man ansieht, rechts, was man tut; Spaltengrenze überall an derselben Stelle (gemeinsam.md „Satz“)
% v0.17: Fußeintrag „Nr: Wert“ – kurz in einer Box (nie getrennt), lang umbrechbar mit der Nummer am Anfang
% (vorher \mbox: lange Lösungen liefen über den Rand, AS40 138 pt); robust, weil \fusshilfe in die .aux schreibt
\NewDocumentCommand{\gfuss}{m m}{\sbox0{#1:~#2}\ifdim\wd0<0.45\textwidth\mbox{#1:~#2}\else\mbox{#1:}\nobreak\ #2\fi}
\newsavebox{\gbox}\newlength{\gLinks}
% v0.10 (Lehrer 13.10.): keine starre Halbierung – die linke Spalte ist so breit wie die Figur (+ 6 mm), die
% Antwort steht direkt daneben; ist die Figur breiter als 0,62 Zeile, steht die Antwort darunter
\newcommand{\gzwei}[2]{\par\smallskip\noindent\sbox{\gbox}{#1}%
  \setlength{\gLinks}{\dimexpr\wd\gbox+6mm\relax}%
  \ifdim\gLinks>0.62\linewidth
    \usebox{\gbox}\par\smallskip\noindent\begin{minipage}{\linewidth}#2\end{minipage}%
  \else
    \begin{minipage}[t]{\gLinks}\vspace{0pt}\usebox{\gbox}\end{minipage}%
    \begin{minipage}[t]{\dimexpr\linewidth-\gLinks\relax}\vspace{0pt}\raggedright #2\end{minipage}%
  \fi\par}
% 16.10. (Lehrer, PA10 S. 1): Text links, Koordinatensystem rechts daneben; ist das System breiter als 0,58 Zeile,
% steht es wie bisher unter dem Text
% Passt das System knapp nicht mehr auf die Seite (bis 25 % zu hoch), wird es auf den Rest verkleinert statt
% die Seite halb leer zu lassen
\newlength{\gRest}\newlength{\gHoch}
\newcommand{\gtextks}[2]{\sbox{\gbox}{#1}%
  \setlength{\gHoch}{\dimexpr\ht\gbox+\dp\gbox\relax}%
  \ifdim\pagegoal<\maxdimen\setlength{\gRest}{\dimexpr\pagegoal-\pagetotal-4\baselineskip\relax}%
    \ifdim\gHoch>\gRest\ifdim\gRest>0.75\gHoch
      \sbox{\gbox}{\resizebox*{!}{\gRest}{\usebox{\gbox}}}\fi\fi\fi
  \ifdim\wd\gbox>0.58\linewidth #2\par\smallskip\noindent\gpassend{\usebox{\gbox}}\par
  \else\begin{minipage}[t]{\dimexpr\linewidth-\wd\gbox-5mm\relax}\raggedright #2\end{minipage}\hfill
    \begin{minipage}[t]{\wd\gbox}\vspace{0pt}\usebox{\gbox}\end{minipage}\par\fi}
% v0.9 (12.10., pruefung.md § 5 „Marke mit B, Bedingungen rechts“): drittes Argument von pfaufg = Bedingungen
% der Prüfung (v0.10: \pfohnetr = TR durchgestrichen, „FOR“), klein grau im rechten Rand auf der Grundlinie der Nummer, Zeilen mit \\
\newcommand{\gbedrand}[2]{\ifblank{#2}{}{\rlap{\hspace*{\dimexpr#1+2mm\relax}%
  \parbox[t]{15mm}{\raggedright\scriptsize\leavevmode\color{mbgrau}\strut#2}}}}
% Teil eines Bündels (a, b …): Marke links im Rand, Bedingungen rechts im Rand, auf der Zeile des Teils
\newcommand{\gteilrand}[2]{\llap{\mbpfrandmarke{#1}\hspace*{\mbpfnr}}\gbedrand{\linewidth}{#2}}
% Teil eines Bündels: Buchstabe in eigener Spalte, Inhalt eingerückt (\gEin = Einzug, damit die Spaltengrenze von
% \gzwei an derselben Stelle bleibt)
\newenvironment{gbteil}[3]{\par\smallskip\noindent\gteilrand{#2}{#3}\makebox[\mbpfteil][l]{\textbf{#1}}%
  \begin{minipage}[t]{\dimexpr\linewidth-\mbpfteil\relax}\setlength{\gEin}{\mbpfteil}\raggedright\strut\ignorespaces}%
 {\end{minipage}\par\addvspace{6pt}}% v0.18 (PG80 Nr. 1): Abstand, damit die Linie von a) nicht an b) rutscht
% Grafik höchstens so breit wie die Spalte (Bündel nebeneinander, Koordinatensystem ohne Karo)
\newsavebox{\gpbox}
\newcommand{\gpassend}[1]{\sbox{\gpbox}{#1}\ifdim\wd\gpbox>\linewidth\typeout{GPASSEND: verkleinert \the\wd\gpbox\space auf \the\linewidth}\resizebox{\linewidth}{!}{\usebox{\gpbox}}\else\usebox{\gpbox}\fi}
% Skizzenfeld (v0.10, Lehrer 13.10.): Karo ab dem linken Rand über die Zeile, „Skizze“ klein grau oben links im Feld
% (wie \pfskizzenfeld der Vorlage); der rechte Rand bleibt den Bedingungen
\newcommand{\gskizze}[1]{\par\smallskip\noindent\begin{tikzpicture}
  \clip (0,0) rectangle ({\linewidth-1mm},{#1*10mm});
  \draw[black!16,line width=0.3pt,step=5mm] (0,0) grid ({\linewidth},{#1*10mm});
  \node[anchor=north west,inner sep=2pt,font=\scriptsize,text=mbgrau,fill=white] at (0,{#1*10mm}) {Skizze};
  \end{tikzpicture}\par}
\RenewDocumentEnvironment{pfaufg}{m m m}{%
  \begin{lrbox}{\mbaufgabebox}\begin{minipage}[t]{\linewidth}%
  \gSchrift\noindent\hspace*{\gEin}\mbpfrandmarke{#1}\gbedrand{\linewidth-\gEin}{#3}%
  \begin{minipage}[t]{\mbpfnr}\strut\textbf{#2}\end{minipage}%
  \begin{minipage}[t]{\dimexpr\linewidth-\gEin-\mbpfnr\relax}\raggedright\strut\ignorespaces}%
 {\end{minipage}%
  \end{minipage}\end{lrbox}%
  \setlength{\mbaufgabehoehe}{\dimexpr\ht\mbaufgabebox+\dp\mbaufgabebox\relax}%
  \ifdefempty{\gkopfnach}{\par\addvspace{5pt}\Needspace*{\mbaufgabehoehe}}{%
    % Stufenkopf, der mit dieser (ersten) Aufgabe zusammenbleibt: Platz für beide verlangen, dann den Kopf setzen
    % (12.10., Quadratische: ein Kopf stand allein unten, die Aufgabe mit Koordinatensystem auf der nächsten Seite)
    % v0.17 (AS40 „Terme deuten“ allein unten): hinter dem Kopf steht das Sprungziel (\special, ein Knoten, nach
    % dem die Seite brechen durfte) – darum \nobreak nach dem Kopf; Vorhalt 56 pt = Kopf mit Linie und Abständen
    \par\Needspace*{\dimexpr\mbaufgabehoehe+56pt\relax}\gkopfnach\par\nobreak\gdef\gkopfnach{}\addvspace{5pt}}%
  \noindent\usebox{\mbaufgabebox}\par\addvspace{5pt}}
% v0.18 (plan.md § 10, 15.10. l): Abschnittsbuchstabe A, B, C vor dem Namen im Kopf – als Vorsatz, damit der
% Name im TeX gleich bleibt (Bestellkarte, Sprungziele lesen ihn)
\def\gbstakt{}\newcommand{\gbstvor}[1]{\gdef\gbstakt{#1}}
\NewCommandCopy\gpfstufekopfalt\pfstufekopf
\RenewDocumentCommand{\pfstufekopf}{o m m}{\IfValueTF{#1}{\gpfstufekopfalt[#1]}{\gpfstufekopfalt}%
  {\ifx\gbstakt\empty\else\gbstakt\enspace\fi #2}{#3}\gdef\gbstakt{}}
% v0.18 (16.10. e): Begriff grau, ohne Nummer, ohne Rahmen über der ersten Aufgabe der Stufe
\newcommand{\gbegriff}[1]{\par\nobreak\noindent\hspace*{\mbpfnr}\parbox[t]{\dimexpr\linewidth-\mbpfnr\relax}{\raggedright\small\color{mbgrau}#1}\par\nobreak\addvspace{3pt}}
% v0.18 (16.10. g): Lösungsblatt – Kopf mit kräftigem Strich, Abschnitt klein grau als Zwischenzeile
\newcommand{\glkopf}[1]{\noindent{\large\bfseries Lösungen \textperiodcentered{} #1}\par\vspace{3pt}\hrule height 1.2pt\par\medskip}
\newcommand{\glzwischen}[1]{\par\addvspace{8pt}\Needspace*{4\baselineskip}\noindent{\footnotesize\color{mbgrau}#1}\par\nobreak\vspace{1pt}}
% v0.18 (15.10. o, w): Blattende – Kern ohne Rahmen, dann graue Zeile „Mehr: n Aufgaben“
\newcommand{\gkern}[1]{\par\addvspace{14pt}\Needspace*{6\baselineskip}\noindent{\bfseries Kern}\par\nobreak\smallskip
  {\small\raggedright #1\par}}
\newcommand{\gmehr}[1]{\par\addvspace{10pt}\noindent{\footnotesize\color{mbgrau}Mehr: #1 Aufgaben}\par}
% Stufenkopf vormerken: \gkopfvor{<Kopf>} – die nächste pfaufg setzt ihn, sobald beide zusammen Platz haben
\def\gkopfnach{}
\newcommand{\gkopfvor}[1]{\gdef\gkopfnach{#1}}
% Teilaufgabe mit Skizze links in gzwei: Buchstabe in eigener Spalte, darunter Text und Skizze
\newcommand{\gteil}[2]{\makebox[\mbpfteil][l]{\textbf{#1}}\parbox[t]{\dimexpr0.5\textwidth-\gEin-\mbpfnr-5mm-\mbpfteil\relax}{\raggedright #2}}
% Stufenkopf der Vorlage (mathblatt.sty 2026-10-12: Linie, Abstand, 13 pt, \pfjahre) ohne ihren Vorhalt von
% 14 Zeilen: der Setzer fordert vor dem Kopf 0,16 Seitenhöhe (Kopf + erste Aufgabe), mehr nicht
\setlength{\mbpfstufevorhalt}{0pt}
\makeatother
\newcommand{\karo}[1]{\par\smallskip\noindent\begin{tikzpicture}
  \clip (0,0) rectangle ({\linewidth-1mm},{#1*10mm});
  \draw[black!16,line width=0.3pt,step=5mm] (0,0) grid ({\linewidth},{#1*10mm});
  \end{tikzpicture}\par}
\newcommand{\kk}[1]{$\square$\,#1\hspace{0.9em}}
\newcommand{\linie}[1]{\,{\color{black!45}\rule[-1pt]{#1}{0.4pt}}\,}
\newcommand{\pkt}{\textperiodcentered{}}
% v0.14 (14.10.): Stufenkopf mit fünf Punkten (Option = Jahre für die Bestellkarte) oder „zuletzt 19“
\NewDocumentCommand{\gpunkte}{o m}{#2}
\newcommand{\gzuletzt}[1]{zuletzt #1}
% v0.7 (12.10.): Schülerfrage unter dem Stufenkopf; Ankreuzoption als Bild (Buchstabe und Kästchen darunter)
\newcommand{\gfrage}[1]{\par\nobreak\noindent\hspace*{\mbpfnr}{\small\itshape #1}\par\nobreak\addvspace{2pt}}
\newlength{\gbildbreite}\setlength{\gbildbreite}{0.32\linewidth}% Setzer setzt sie je Reihe (vier Bilder: 0,24)
\newcommand{\gbildoption}[3]{\begin{minipage}[t][#3][b]{\gbildbreite}\centering #2\par\smallskip\mbox{#1\,$\square$}\end{minipage}}
% Lineare (12.10.): Bild mit Buchstabe ohne Kästchen (Diagramm zuordnen); Aussage mit „wahr / falsch“ zum Ankreuzen
\newcommand{\gbildname}[3]{\begin{minipage}[t][#3][b]{\gbildbreite}\centering #2\par\smallskip\mbox{\textbf{#1}}\end{minipage}}
\newcommand{\gwahrfalsch}[2]{\par\noindent\hangindent2.2em\hangafter1\makebox[2.2em][l]{(#1)}#2\hfill\mbox{wahr $\square$\quad falsch $\square$}\par}

\newcommand{\gunter}[2]{\par\addvspace{6pt}\noindent{\bfseries #1\enspace #2}\par\nobreak\vspace{2pt}}
\newcommand{\gsternzeile}{\par\vspace{-4pt}\noindent{\footnotesize\color{mbgrau}* nur FOR}\par\medskip}
% 17.10.: Skizze oben bündig mit der ersten Textzeile (beide Spalten hängen an der Oberkante); der Rechenraum (Karo)
% beginnt direkt unter dem Text, links neben dem Graph, und läuft unter dem Graph über die ganze Zeile weiter
% (Streifen links = auf 5 mm aufgerundete Höhe des Graphen unter dem Text, Raster durchgehend; die Fläche des Streifens
% wird von der Karohöhe unter dem Graph abgezogen, auf 0,5 cm aufgerundet – die Gesamtfläche Rechenraum bleibt mindestens #3 × Zeile); #3 = Zeilen Karo in cm
% unter dem Graph (0 = kein Karo). Kein Verkleinern des Graphen.
\newsavebox{\gtbox}\newlength{\gWl}\newlength{\gDh}
\newcommand{\gkaroohne}[1]{\nointerlineskip\noindent\begin{tikzpicture}
  \clip (0,0) rectangle ({\linewidth-1mm},{#1*10mm});
  \draw[black!16,line width=0.3pt,step=5mm] (0,0) grid ({\linewidth},{#1*10mm});
  \end{tikzpicture}\par}
\newcommand{\gtextkaro}[3]{\sbox{\gbox}{#1}%
  \ifdim\wd\gbox>0.58\linewidth #2\par\smallskip\noindent\gpassend{\usebox{\gbox}}\par\ifnum#3>0 \karo{#3}\fi
  \else
    \setlength{\gWl}{\dimexpr\linewidth-\wd\gbox-5mm\relax}%
    \sbox{\gtbox}{\begin{minipage}[t]{\gWl}\raggedright #2\par\end{minipage}}%
    \setlength{\gDh}{\dimexpr\ht\gbox+\dp\gbox-\ht\gtbox-\dp\gtbox-6pt\relax}%
    \leavevmode
    \ifnum#3>0\ifdim\gDh>0pt
      \pgfmathsetmacro{\gxs}{floor(\gWl/5mm)*5mm}\pgfmathsetmacro{\gys}{ceil(\gDh/5mm)*5mm}%
      \begin{minipage}[t]{\gWl}\usebox{\gtbox}\par\vspace{6pt}\nointerlineskip\noindent
        \begin{tikzpicture}
        \clip (0,0) rectangle (\gxs pt,\gys pt);
        \draw[black!16,line width=0.3pt,step=5mm] (0,0) grid (\gxs pt,\gys pt);
        \end{tikzpicture}\par\end{minipage}%
    \else\usebox{\gtbox}\fi\else\usebox{\gtbox}\fi
    \hfill\begin{minipage}[t]{\wd\gbox}\vspace{-\dimexpr\ht\strutbox+2.4ex\relax}\usebox{\gbox}\end{minipage}\par
    \ifnum#3>0 \ifdim\gDh>0pt
      \pgfmathsetmacro{\gnn}{max(0,ceil(2*(#3-(\gys/28.4528)*(\gxs/\linewidth)))/2)}%
      \ifdim\gnn pt>0pt\gkaroohne{\gnn}\fi
    \else\karo{#3}\fi\fi
  \fi}
\begin{document}

\pfheftstil
\pflaufkopf[P10 \pkt{} Prüfung]{Exponentielles Wachstum}{PA50}\pfseitevon
\pfheftkopf{Exponentielles Wachstum}{}
\pfstufenzeile{Punkte und Funktionswerte \pkt{} Exponentielles Wachstum \pkt{} Exponentialfunktion \pkt{} Wachstumsmodelle}
\gsternzeile

% ===================== A Punkte und Funktionswerte =====================
\gkopfvor{\gbstvor{A}\pfstufekopf{Punkte und Funktionswerte}{}}
% DRUCK 2026-FOR-K7a deckt 2026-FOR-K7a 2026-EBR-K7a
\begin{pfaufg}{nach~P10\\’26}{1.}{}
\gtextkaro{{\renewcommand{\arraystretch}{1.25}\begin{tabular}{|l|c|c|c|c|}\hline Jahr & 2026 & 2027 & 2028 & 2029\\\hline Miete in € & \leerzelle & 662,35 & 674,93 & \leerzelle\\\hline\end{tabular}}}{Can zieht Anfang 2026 in eine neue Wohnung. Er erhält einen Mietvertrag, bei dem sich die Miete jedes Jahr um 1,9 \% erhöht. Im ersten Jahr zahlt Can 650 € monatlich. Vervollständige die Tabelle.}{4}
\fusshilfe{\gfuss{1}{2026: 650,00 €; 2029: 687,75 €}}
\end{pfaufg}
% DRUCK 2025-OS-K7a deckt 2025-OS-K7a
\begin{pfaufg}{nach~P10\\’25}{2.}{}
\gtextkaro{\hbox{\vtop{\vspace{0pt}\hbox{{\renewcommand{\arraystretch}{1.25}\begin{tabular}{|l|c|c|c|c|c|}\hline Zeit in Stunden & 0 & 1 & 2 & 3 & 4\\\hline Anzahl der Bakterien & 80 & 112 & 157 & 220 & 308\\\hline\end{tabular}}}}\hspace{1.2cm}\vtop{\vspace{0pt}\hbox{\begin{ksys}[xmin=0,xmax=5,ymin=0,ymax=400,xstep=1,ystep=25,karo=0.5,xlabel={Zeit in h},ylabel={Anzahl der Bakterien}]\end{ksys}}}}}{Das exponentielle Wachstum eines Bakterienstammes wird beobachtet. Die Tabelle zeigt das Wachstum des Bakterienstammes in den ersten 4 Stunden. Stelle das Bakterienwachstum im Koordinatensystem dar.}{0}
\fusshilfe{\gfuss{2}{Punkte (0|80), (1|112), (2|157), (3|220), (4|308)}}
\end{pfaufg}
\par\medskip\noindent\textcolor{mbgrau}{\rule{\linewidth}{0.3pt}}\par\nobreak\noindent{\footnotesize\textcolor{mbgrau}{2016–2021}}\par\nobreak\smallskip
% DRUCK 2020-OS-K4a deckt 2020-OS-K4a
\begin{pfaufg}{nach~P10\\’20}{3.}{}
\gtextkaro{{\renewcommand{\arraystretch}{1.25}\begin{tabular}{|>{\raggedright\arraybackslash}p{5.2cm}|c|c|c|c|}\hline Jahr & 2019 & 2020 & 2021 & 2022\\\hline Jährlicher Wasserverbrauch pro Person in Tausend Litern & \leerzelle & 55,6 & 57,3 & \leerzelle\\\hline\end{tabular}}}{Eine Studie aus dem Jahr 2019 besagt, dass der weltweite Wasserverbrauch pro Person in Zukunft um 3 \% pro Jahr steigen wird. Im Jahr 2019 wurde ein jährlicher Wasserverbrauch von durchschnittlich 54 Tausend Liter pro Person ermittelt. Ergänze die Tabelle entsprechend der Angaben aus dem Text.}{4}
\fusshilfe{\gfuss{3}{2019: 54; 2022: 59,0}}
\end{pfaufg}
% DRUCK 2019-OS-K7a deckt 2019-OS-K7a
\begin{pfaufg}{nach~P10\\’19}{4.}{}
\gtextkaro{{\renewcommand{\arraystretch}{1.25}\begin{tabular}{|l|c|c|c|c|c|}\hline Zeit in Wochen & 0 & 1 & 2 & \dots & 5\\\hline Masse in kg & \leerzelle & 10,400 & \leerzelle & – & \leerzelle\\\hline\end{tabular}}}{Ein Ferkel hat eine Masse von 10 kg. Erfahrungsgemäß nimmt die Masse des Ferkels wöchentlich um ca. 4 \% zu. Am Anfang jeder Woche wird die Masse des Ferkels kontrolliert. Ergänze die drei Werte in der Wertetabelle.}{5}
\fusshilfe{\gfuss{4}{Woche 0: 10,000 kg; Woche 2: 10,816 kg; 12,167 kg}}
\end{pfaufg}
% DRUCK 2017-OS-K7a deckt 2017-OS-K7a
% DRUCK 2017-OS-K7d deckt 2017-OS-K7d
\begin{pfaufg}{nach~P10\\’17}{5.}{}
\gtextkaro{{\renewcommand{\arraystretch}{1.25}\begin{tabular}{|l|c|c|c|c|c|c|c|}\hline Höhe in km & 0 & 1 & 2 & 3 & 5 & 8 & 10\\\hline Luftdruck in hPa & 1000 & 870 & \leerzelle & 659 & 498 & \leerzelle & 248\\\hline\end{tabular}}}{In Höhe des Meeresspiegels (0 km Höhe) beträgt der Luftdruck ungefähr 1000 hPa (Hektopascal). Je höher man in der Erdatmosphäre steigt, desto geringer wird der Luftdruck. Er nimmt pro Kilometer um 13 \% ab.\begin{gbteil}{a)}{}{}Ergänze die fehlenden Werte in der Tabelle.\end{gbteil}}{4}
\begin{gbteil}{b)}{}{}
Ein Bergsteiger misst einen Luftdruck von 573 hPa. Entscheide, ob er sich in ca. 4 km Höhe befinden kann. Begründe.
\karo{4}
\end{gbteil}
\fusshilfe{\gfuss{5}{a) 757 hPa; 328 hPa; b) ja}}
\end{pfaufg}
% DRUCK 2016-OS-K4e deckt 2016-OS-K4e
\begin{pfaufg}{nach~P10\\’16}{6.\textsuperscript{*}}{}
Mit Hilfe radioaktiver Strahlung können Erkrankungen der Schilddrüse erkannt werden. Dazu wird dem Patienten eine Flüssigkeit mit einer kleinen Menge des radioaktiven Elements Technetium gespritzt. Pro Stunde nimmt die Masse des Technetiums um ca. 11 \% ab. Zu Beginn (0 h) sind es 5,00 mg. Gib an, wie viel Milligramm Technetium nach einem Tag (24 h) noch vorhanden sind.
\karo{3}
\par\medskip\noindent\hfill m $=$\linie{45mm}
\fusshilfe{\gfuss{6}{0,31 mg}}
\end{pfaufg}

% ===================== B Exponentielles Wachstum =====================
\gkopfvor{\gbstvor{B}\pfstufekopf{Exponentielles Wachstum}{}}
% DRUCK 2018-OS-K2a deckt 2018-OS-K2a
% DRUCK 2018-OS-K2b deckt 2018-OS-K2b
\begin{pfaufg}{nach~P10\\’18}{7.}{}
\gtextkaro{{\renewcommand{\arraystretch}{1.25}\begin{tabular}{|l|r|r|r|r|r|}\hline Jahr & 2013 & 2014 & 2015 & 2016 & 2017\\\hline Umsatz in € & 600\,000 & 648\,000 & 699\,840 & 755\,827 & 816\,293\\\hline\end{tabular}}}{In der Tabelle stellt ein Berliner Start-up-Unternehmen die Entwicklung seiner Umsätze dar. Das Unternehmen gibt an, dass es immer ein jährliches Wachstum von 8 \% erreicht.\begin{gbteil}{a)}{}{}Zeige, dass diese Aussage für das Jahr 2014 richtig ist.\end{gbteil}}{3}
\begin{gbteil}{b)}{}{}
Eine Überprüfung ergab, dass die Umsätze tatsächlich jährlich um 8 \% steigen. Entscheide und begründe, ob lineares oder exponentielles Wachstum vorliegt.
\karo{4}
\end{gbteil}
\fusshilfe{\gfuss{7}{a) stimmt; b) exponentiell}}
\end{pfaufg}\vspace{-6pt}
% DRUCK 2017-OS-K7c deckt 2017-OS-K7c
\begin{pfaufg}{nach~P10\\’17}{8.\textsuperscript{*}}{}
In Höhe des Meeresspiegels (0 km Höhe) beträgt der Luftdruck ungefähr 1000 hPa (Hektopascal). Je höher man in der Erdatmosphäre steigt, desto geringer wird der Luftdruck. Er nimmt pro Kilometer um 13 \% ab. Die Abnahme des Luftdrucks kann mit einer Funktionsgleichung beschrieben werden. Kreuze die richtige Gleichung an:\pfkreuzab{}{$y = 0{,}87 \cdot x$}\pfkreuzab{}{$y = 1000 \cdot 0{,}87^x$}\pfkreuzab{}{$y = x^{1{,}13}$}\pfkreuzab{}{$y = 1000 \cdot 1{,}13^x$}
\fusshilfe{\gfuss{8}{$y = 1000 \cdot 0{,}87^x$}}
\end{pfaufg}\vspace{-6pt}
% DRUCK 2016-OS-K4d deckt 2016-OS-K4d
\begin{pfaufg}{nach~P10\\’16}{9.\textsuperscript{*}}{}
Mit Hilfe radioaktiver Strahlung können Erkrankungen der Schilddrüse erkannt werden. Dazu wird dem Patienten eine Flüssigkeit mit einer kleinen Menge des radioaktiven Elements Technetium gespritzt. Pro Stunde nimmt die Masse des Technetiums um ca. 11 \% ab. Entscheide, welche der beiden Aussagen über den Abnahmeprozess richtig ist. Kreuze an:\pfkreuzab{}{Es handelt sich um eine lineare Abnahme}\pfkreuzab{}{Es handelt sich um eine exponentielle Abnahme}
\par\smallskip\noindent Begründe deine Entscheidung.
\karo{3}
\fusshilfe{\gfuss{9}{exponentielle Abnahme}}
\end{pfaufg}\vspace{-6pt}
% DRUCK 2019-OS-K7b deckt 2019-OS-K7b
\begin{pfaufg}{nach~P10\\’19}{10.\textsuperscript{*}}{}
\gtextkaro{{\renewcommand{\arraystretch}{1.6}\begin{tabular}{|c|p{10cm}|}\hline & \textbf{Bedeutung}\\\hline $10$ & \\\hline $1{,}04$ & \\\hline $x$ & \\\hline\end{tabular}}}{Ein Ferkel hat eine Masse von 10 kg. Erfahrungsgemäß nimmt die Masse des Ferkels wöchentlich um ca. 4 \% zu. Um die Masse auszurechnen, die das Ferkel nach mehreren Wochen hat, stellt Landwirtin Wilma die Funktionsgleichung $f(x) = 10 \cdot 1{,}04^x$ auf. Gib die Bedeutung der Bestandteile der Funktionsgleichung an.}{0}
\fusshilfe{\gfuss{10}{10 – Anfangsmasse in kg; 1,04 – Wachstumsfaktor, +4 \% je Woche; x – Zeit in Wochen}}
\end{pfaufg}\vspace{-6pt}

% ===================== C Exponentialfunktion =====================
\gkopfvor{\gbstvor{C}\pfstufekopf{Exponentialfunktion}{}}
% DRUCK 2026-FOR-K7b deckt 2026-FOR-K7b 2026-EBR-K7b
\begin{pfaufg}{nach~P10\\’26}{11.}{}
Can zieht Anfang 2026 in eine neue Wohnung. Er erhält einen Mietvertrag, bei dem sich die Miete jedes Jahr um 1,9 \% erhöht. Im ersten Jahr zahlt Can 650 € monatlich. Entscheide, welcher Graph zur Entwicklung der Miete passt. Kreuze den passenden Graphen an. Begründe für die beiden anderen Graphen, warum sie nicht zum Sachverhalt passen.
\par\smallskip\noindent\setlength{\gbildbreite}{\dimexpr0.96\linewidth/3\relax}\sbox{\gbox}{\vbox{\begin{tikzpicture}[x=0.45cm,y=0.45cm,line width=0.5pt]\node[font=\small\bfseries] at (4,6.9) {A};\draw[black!20,line width=0.3pt,step=1] (0,0) grid (8,6);\draw[-{Stealth[length=1.5mm]}] (0,0)--(8.5,0);\draw[-{Stealth[length=1.5mm]}] (0,0)--(0,6.5);\node[below,font=\tiny] at (4.6,0) {Zeit in Jahren};\node[right,font=\tiny] at (0,6.4) {Miete in €};\draw[line width=0.8pt,domain=0:7.6,samples=40,smooth] plot (\x,{0.25*(exp(0.43*\x)-1)});\end{tikzpicture}\par\begin{tikzpicture}[x=0.45cm,y=0.45cm,line width=0.5pt]\node[font=\small\bfseries] at (4,6.9) {B};\draw[black!20,line width=0.3pt,step=1] (0,0) grid (8,6);\draw[-{Stealth[length=1.5mm]}] (0,0)--(8.5,0);\draw[-{Stealth[length=1.5mm]}] (0,0)--(0,6.5);\node[below,font=\tiny] at (4.6,0) {Zeit in Jahren};\node[right,font=\tiny] at (0,6.4) {Miete in €};\draw[line width=0.8pt,domain=0:7.6,samples=40,smooth] plot (\x,{1.2*exp(0.2*\x)});\end{tikzpicture}\par\begin{tikzpicture}[x=0.45cm,y=0.45cm,line width=0.5pt]\node[font=\small\bfseries] at (4,6.9) {C};\draw[black!20,line width=0.3pt,step=1] (0,0) grid (8,6);\draw[-{Stealth[length=1.5mm]}] (0,0)--(8.5,0);\draw[-{Stealth[length=1.5mm]}] (0,0)--(0,6.5);\node[below,font=\tiny] at (4.6,0) {Zeit in Jahren};\node[right,font=\tiny] at (0,6.4) {Miete in €};\draw[line width=0.8pt] (0,1.3)--(7.6,5.2);\end{tikzpicture}}}\gbildoption{}{\begin{tikzpicture}[x=0.45cm,y=0.45cm,line width=0.5pt]\node[font=\small\bfseries] at (4,6.9) {A};\draw[black!20,line width=0.3pt,step=1] (0,0) grid (8,6);\draw[-{Stealth[length=1.5mm]}] (0,0)--(8.5,0);\draw[-{Stealth[length=1.5mm]}] (0,0)--(0,6.5);\node[below,font=\tiny] at (4.6,0) {Zeit in Jahren};\node[right,font=\tiny] at (0,6.4) {Miete in €};\draw[line width=0.8pt,domain=0:7.6,samples=40,smooth] plot (\x,{0.25*(exp(0.43*\x)-1)});\end{tikzpicture}}{\dimexpr\ht\gbox/3+\baselineskip\relax}\hfill\gbildoption{}{\begin{tikzpicture}[x=0.45cm,y=0.45cm,line width=0.5pt]\node[font=\small\bfseries] at (4,6.9) {B};\draw[black!20,line width=0.3pt,step=1] (0,0) grid (8,6);\draw[-{Stealth[length=1.5mm]}] (0,0)--(8.5,0);\draw[-{Stealth[length=1.5mm]}] (0,0)--(0,6.5);\node[below,font=\tiny] at (4.6,0) {Zeit in Jahren};\node[right,font=\tiny] at (0,6.4) {Miete in €};\draw[line width=0.8pt,domain=0:7.6,samples=40,smooth] plot (\x,{1.2*exp(0.2*\x)});\end{tikzpicture}}{\dimexpr\ht\gbox/3+\baselineskip\relax}\hfill\gbildoption{}{\begin{tikzpicture}[x=0.45cm,y=0.45cm,line width=0.5pt]\node[font=\small\bfseries] at (4,6.9) {C};\draw[black!20,line width=0.3pt,step=1] (0,0) grid (8,6);\draw[-{Stealth[length=1.5mm]}] (0,0)--(8.5,0);\draw[-{Stealth[length=1.5mm]}] (0,0)--(0,6.5);\node[below,font=\tiny] at (4.6,0) {Zeit in Jahren};\node[right,font=\tiny] at (0,6.4) {Miete in €};\draw[line width=0.8pt] (0,1.3)--(7.6,5.2);\end{tikzpicture}}{\dimexpr\ht\gbox/3+\baselineskip\relax}\par
\karo{4}
\fusshilfe{\gfuss{11}{Graph B; Graph A – Start bei 0 € statt 650 €; Graph C – linear statt exponentiell}}
\end{pfaufg}
% DRUCK 2019-OS-K7c deckt 2019-OS-K7c
\begin{pfaufg}{nach~P10\\’19}{12.\textsuperscript{*}}{}
\gtextkaro{{\renewcommand{\arraystretch}{1.3}\begin{tabular}{|m{4.1cm}|>{\centering\arraybackslash}m{1.1cm}|>{\centering\arraybackslash}m{1.5cm}|m{5.4cm}|}\hline \textbf{Skizze des Graphen} & \textbf{passt} & \textbf{passt nicht} & \textbf{Begründung}\\\hline \vspace{4pt}\begin{tikzpicture}[x=0.45cm,y=0.45cm,line width=0.5pt]\draw[black!20,line width=0.3pt,step=1] (0,0) grid (8,6);\draw[-{Stealth[length=1.5mm]}] (0,0)--(8.5,0);\draw[-{Stealth[length=1.5mm]}] (0,0)--(0,6.5);\node[below,font=\tiny] at (4.6,0) {Zeit in Wochen};\node[right,font=\tiny] at (0,6.4) {Masse in kg};\draw[line width=0.8pt] (0,1.3)--(7.6,5.2);\end{tikzpicture}\vspace{4pt} & $\square$ & $\square$ & \\\hline \vspace{4pt}\begin{tikzpicture}[x=0.45cm,y=0.45cm,line width=0.5pt]\draw[black!20,line width=0.3pt,step=1] (0,0) grid (8,6);\draw[-{Stealth[length=1.5mm]}] (0,0)--(8.5,0);\draw[-{Stealth[length=1.5mm]}] (0,0)--(0,6.5);\node[below,font=\tiny] at (4.6,0) {Zeit in Wochen};\node[right,font=\tiny] at (0,6.4) {Masse in kg};\draw[line width=0.8pt,domain=0:7.6,samples=40,smooth] plot (\x,{1.2*exp(0.2*\x)});\end{tikzpicture}\vspace{4pt} & $\square$ & $\square$ & \\\hline \vspace{4pt}\begin{tikzpicture}[x=0.45cm,y=0.45cm,line width=0.5pt]\draw[black!20,line width=0.3pt,step=1] (0,0) grid (8,6);\draw[-{Stealth[length=1.5mm]}] (0,0)--(8.5,0);\draw[-{Stealth[length=1.5mm]}] (0,0)--(0,6.5);\node[below,font=\tiny] at (4.6,0) {Zeit in Wochen};\node[right,font=\tiny] at (0,6.4) {Masse in kg};\draw[line width=0.8pt,domain=0:7.6,samples=40,smooth] plot (\x,{0.25*(exp(0.43*\x)-1)});\end{tikzpicture}\vspace{4pt} & $\square$ & $\square$ & \\\hline\end{tabular}}}{Ein Ferkel hat eine Masse von 10 kg. Erfahrungsgemäß nimmt die Masse des Ferkels wöchentlich um ca. 4 \% zu. Entscheide für jeden Graphen, ob er zum Wachstum des Ferkels passt. Begründe deine Entscheidung.}{0}
\fusshilfe{\gfuss{12}{Graph 1 – passt nicht; Graph 2 – passt; Graph 3 – passt nicht}}
\end{pfaufg}

% ===================== D Wachstumsmodelle =====================
\gkopfvor{\gbstvor{D}\pfstufekopf{Wachstumsmodelle}{}}
% DRUCK 2026-FOR-K7c deckt 2026-FOR-K7c
\begin{pfaufg}{nach~P10\\’26}{13.\textsuperscript{*}}{}
Can zieht Anfang 2026 in eine neue Wohnung. Er erhält einen Mietvertrag, bei dem sich die Miete jedes Jahr um 1,9 \% erhöht. Im ersten Jahr zahlt Can 650 € monatlich. Stelle eine Funktionsgleichung auf, mit der die monatliche Miete für eine beliebige Anzahl von Jahren berechnet werden kann. Ermittle die monatliche Miete im Jahr 2040.
\karo{4}
\par\medskip\noindent\hfill M(t) $=$\linie{45mm}
\par\medskip\noindent\hfill Miete 2040 $=$\linie{45mm}
\fusshilfe{\gfuss{13}{650 · $1{,}019^{t}$; 845,96 €}}
\end{pfaufg}
% DRUCK 2025-OS-K7b deckt 2025-OS-K7b
\begin{pfaufg}{nach~P10\\’25}{14.\textsuperscript{*}}{}
\gtextkaro{{\renewcommand{\arraystretch}{1.25}\begin{tabular}[t]{|l|c|}\hline Zeit in h & Anzahl\\\hline 0 & 80\\\hline 1 & 112\\\hline 2 & 157\\\hline 3 & 220\\\hline 4 & 308\\\hline\end{tabular}}}{Das exponentielle Wachstum eines Bakterienstammes wird beobachtet. Die Tabelle zeigt das Wachstum des Bakterienstammes in den ersten 4 Stunden. Weise nach, dass der Wachstumsfaktor 1,4 beträgt. Gib eine Gleichung an, die das Bakterienwachstum nach $t$ Stunden beschreibt. Überprüfe rechnerisch, ob folgende Aussage wahr ist: „Nach 24 Stunden ist der Bakterienstamm auf mehr als 250 000 Bakterien angewachsen.“}{5}
\par\medskip\noindent\hfill N(t) $=$\linie{45mm}
\fusshilfe{\gfuss{14}{Faktor 1,4; 80 · $1{,}4^{t}$; wahr}}
\end{pfaufg}
\par\medskip\noindent\textcolor{mbgrau}{\rule{\linewidth}{0.3pt}}\par\nobreak\noindent{\footnotesize\textcolor{mbgrau}{2016–2021}}\par\nobreak\smallskip
% DRUCK 2016-OS-K4c deckt 2016-OS-K4c
\begin{pfaufg}{nach~P10\\’16}{15.}{}
\gtextkaro{{\renewcommand{\arraystretch}{1.25}\begin{tabular}{|l|c|c|c|c|c|c|c|}\hline Zeit in h & 0 & 1 & 2 &  & 5 & 6 & 8\\\hline Masse des Technetiums in mg & 5,00 &  & 3,96 & 3,14 & 2,79 & 2,48 & 1,97\\\hline\end{tabular}}}{Mit Hilfe radioaktiver Strahlung können Erkrankungen der Schilddrüse erkannt werden. Dazu wird dem Patienten eine Flüssigkeit mit einer kleinen Menge des radioaktiven Elements Technetium gespritzt. Pro Stunde nimmt die Masse des Technetiums um ca. 11 \% ab. Die Tabelle zeigt die Abnahme. Gib an, nach welcher Zeit nur noch die Hälfte der Masse des Technetiums im Körper vorhanden ist.}{3}
\par\medskip\noindent\hfill t $=$\linie{45mm}
\fusshilfe{\gfuss{15}{nach etwa 6 h}}
\end{pfaufg}
% DRUCK 2018-OS-K2c deckt 2018-OS-K2c
\begin{pfaufg}{nach~P10\\’18}{16.}{}
\gtextkaro{{\renewcommand{\arraystretch}{1.25}\begin{tabular}{|l|r|r|r|r|r|}\hline Jahr & 2013 & 2014 & 2015 & 2016 & 2017\\\hline Umsatz in € & 600\,000 & 648\,000 & 699\,840 & 755\,827 & 816\,293\\\hline\end{tabular}}}{In der Tabelle stellt ein Berliner Start-up-Unternehmen die Entwicklung seiner Umsätze dar. Die Umsätze steigen jährlich um 8 \%. Gib an, in welchem Jahr der Umsatz erstmals höher als 1\,000\,000 € sein wird (gleiches Wachstum vorausgesetzt). Wie hoch wird in dem ermittelten Jahr der Umsatz sein?}{4}
\par\medskip\noindent\hfill Jahr $=$\linie{25mm}\quad Umsatz $=$\linie{25mm}
\fusshilfe{\gfuss{16}{2020; 1 028 294 €}}
\end{pfaufg}
% DRUCK 2016-OS-K4a deckt 2016-OS-K4a
\begin{pfaufg}{nach~P10\\’16}{17.}{}
\gtextkaro{{\renewcommand{\arraystretch}{1.25}\begin{tabular}{|l|c|c|c|c|c|c|c|}\hline Zeit in h & 0 & 1 & 2 & \leerzelle & 5 & 6 & 8\\\hline Masse des Technetiums in mg & 5,00 & \leerzelle & 3,96 & 3,14 & 2,79 & 2,48 & 1,97\\\hline\end{tabular}}}{Mit Hilfe radioaktiver Strahlung können Erkrankungen der Schilddrüse erkannt werden. Dazu wird dem Patienten eine Flüssigkeit mit einer kleinen Menge des radioaktiven Elements Technetium gespritzt. Pro Stunde nimmt die Masse des Technetiums um ca. 11 \% ab. Vervollständige die Tabelle für diese Abnahme.}{4}
\fusshilfe{\gfuss{17}{4,45 mg; 4 h}}
\end{pfaufg}
% DRUCK 2020-OS-K4b deckt 2020-OS-K4b
\begin{pfaufg}{nach~P10\\’20}{18.}{}
\gtextkaro{\begin{ksys}[xmin=0,xmax=30,ymin=0,ymax=130,xstep=2,ystep=10,karo=0.45,xlabel={Jahre},ylabel={Tausend Liter}]\funktion{54*1.03^\x}{}\end{ksys}}{Eine Studie aus dem Jahr 2019 besagt, dass der weltweite Wasserverbrauch pro Person in Zukunft um 3 \% pro Jahr steigen wird. Im Jahr 2019 wurde ein jährlicher Wasserverbrauch von durchschnittlich 54 Tausend Liter pro Person ermittelt. Der Graph zeigt den weltweiten Wasserverbrauch in Tausend Litern pro Person seit dem Jahr 2019. Ermittle anhand des Graphen, nach wie vielen Jahren sich der Wasserverbrauch pro Person in etwa verdoppelt haben wird. Beschreibe dein Vorgehen.}{4}
\fusshilfe{\gfuss{18}{nach etwa 23 Jahren}}
\end{pfaufg}
% DRUCK 2020-OS-K4c deckt 2020-OS-K4c
\begin{pfaufg}{nach~P10\\’20}{19.\textsuperscript{*}}{}
Eine Studie aus dem Jahr 2019 besagt, dass der weltweite Wasserverbrauch pro Person in Zukunft um 3 \% pro Jahr steigen wird. Die Vorhersage für den weltweiten Wasserverbrauch pro Person kann durch die Gleichung $y = 54 \cdot 1{,}03^x$ beschrieben werden ($x$: Zeit in Jahren, $y$: Wasserverbrauch pro Person in Tausend Litern). Eine andere Studie schätzt, dass der jährliche Wasserverbrauch pro Person von 2019 bis 2033 insgesamt um ca. 40 \% steigen wird. Untersuche, ob beide Studien für das Jahr 2033 den gleichen Wasserverbrauch vorhersagen.
\karo{4}
\fusshilfe{\gfuss{19}{nein}}
\end{pfaufg}
\end{document}